\chapter{Una neurona no es un solo dispositivo}
% versión completa: chapters/18_eetypes.tex

En el primer capítulo clasificamos las neuronas según la sustancia que liberan. Un ingeniero las clasificaría de otra forma: según qué dispositivo hace cada una. Y resulta que bajo un aspecto idéntico se esconde toda una caja de componentes electrónicos.

\section*{Una caja de componentes}

El balde agujereado del segundo capítulo es un acumulador con fugas: suma las señales en una ventana corta. Con una ventana muy pequeña se convierte en un detector de coincidencias. Una neurona que responde solo a los cambios y enseguida se calla es un filtro que deja pasar las variaciones. Una neurona que trabaja sola es un metrónomo, un oscilador; si una entrada ajusta su ritmo, es un oscilador controlado por voltaje. Una neurona que emite ráfagas de clics es un generador de ráfagas que escribe “rayas”. Una neurona que hace clic al compás de la onda sonora es un detector de fase.

\pencilfig{18}{\pencilart{18}}{Bajo un aspecto idéntico se esconde toda una caja de componentes: resistencia, condensador, bobina, oscilador.}

Hay también dispositivos de los que no hablamos antes. El \emph{resonador} se parece a un columpio: responde mejor a empujones de cierto ritmo, como un columpio que solo se balancea si se lo empuja a tiempo. Dentro de la neurona, el papel del resorte lo cumple una corriente lenta que se opone a los cambios. El \emph{acumulador sin fugas} mantiene un valor mucho tiempo, como un condensador sin fugas; así funciona la red que sostiene la posición de los ojos cuando usted mira hacia un lado. Pero para eso la retroalimentación tiene que estar ajustada casi a la perfección, y es un ajuste caro. El \emph{cerrojo} es una neurona a la que un empujón breve pasa a un trabajo prolongado hasta que la apagan, como un interruptor con enclavamiento; en las neuronas que controlan los músculos, ese modo lo activa la serotonina.

\section*{Un dispositivo camaleón}

El descubrimiento principal de esta clasificación: no son piezas distintas, sino distintos \emph{modos} de la misma pieza. Una misma neurona puede ser acumulador de día y generador de ráfagas de noche, según lo que le digan las emisoras de radio. Un ingeniero llamaría al cerebro una matriz lógica reprogramable: el circuito es uno, y los ajustes cambian la función de las piezas.

Que las neuronas saben trabajar en todos estos modos está medido. Cómo usa el cerebro el cambio de modos para calcular es, en su mayor parte, una hipótesis.

\fullversion{18}
